\section{Conductor doubling and rigidity of the new dyadic layer}
\label{sec:dyadic}

This section contains the new structural input. Doubling an odd conductor
does not enlarge its cyclotomic field, whereas doubling an even conductor
does. The two cases must be treated differently.

\subsection{Odd-conductor doubling}

\begin{theorem}[Exact conductor-doubling identity]\label{thm:doubling}
Let $m$ be odd and $(a,2m)=1$. In
$\Lambda^2\A(\Q(\zeta_m))$,
\begin{equation}\label{eq:doubling}
 \beta_{2m}(a)=3\beta_m(a)-\beta_m(2a)-\beta_m(a/2),
\end{equation}
where $a/2$ denotes multiplication by the inverse of $2$ modulo $m$.
The formula also holds for $m=1$ under the zero-symbol convention.
\end{theorem}
\begin{proof}
For $m>1$ split $\mu_{2m}$ into $\mu_m$ and $-\mu_m$. The first
part of the sum defining $\beta_{2m}(a)$ is $\beta_m(a)$. Set
$u_k=\cl{1-\zeta_m^k}$ for $k\not\equiv0\pmod m$. For these $k$,
\[
 \cl{1+\zeta_m^k}=u_{2k}-u_k.
\]
Since $a$ is odd, the second part, apart from the zero wedge at the root
$-1$, is
\[
 \sum_{k=1}^{m-1}(u_{2ak}-u_{ak})\wedge(u_{2k}-u_k).
\]
The two diagonal terms each give $\beta_m(a)$ by reindexing. The cross
terms give $-\beta_m(2a)$ and $-\beta_m(a/2)$. Including the first
part proves \eqref{eq:doubling}. No logarithm branch is used in this
multiplicative-group calculation.
\end{proof}

\subsection{Even-conductor descent is rigid}

\begin{theorem}[New dyadic layer rigidity]\label{thm:layer}
Let $m$ be even and $(a,2m)=1$. Then
\begin{equation}\label{eq:layer}
 \beta_{2m}(a)\in\Lambda^2\A(\Q(\zeta_m))
 \quad\Longleftrightarrow\quad
 \beta_{2m}(a)=0
 \quad\Longleftrightarrow\quad
 a^2\equiv\pm1\pmod{2m}.
\end{equation}
\end{theorem}
\begin{proof}
For $m=2$, all the symbols at conductor $4$ vanish, so the assertion
is immediate. Suppose $m\ge4$. The natural map
$G_{2m}\to G_m$ has kernel $\{1,h\}$, where $h$ is the class of
$1+m$. This is a nontrivial element of order $2$.

In the real regular representation $\R^{G_{2m}}$, logarithm vectors
coming from $\Q(\zeta_m)$ lie in the $+1$ eigenspace of $P_h$.
Consequently, a skew matrix arising from an exterior product of such
vectors annihilates the $-1$ eigenspace. If $\beta_{2m}(a)$ descends,
\eqref{eq:skew-image} therefore gives
\[
 A_{2m}(P_a-P_{a\inv})(I-P_h)=0.
\]
The Gram matrix is invertible. The faithful regular representation of
the group algebra then gives
\begin{equation}\label{eq:four-terms}
 (\ee_a-\ee_{a\inv})(\ee_1-\ee_h)=0.
\end{equation}
If $a\ne a\inv$ in $G_{2m}$, the coefficient of $\ee_a$ in the left
side is positive: the negative term $-\ee_{a\inv}$ is distinct,
$-\ee_{ah}$ is distinct because $h\ne1$, and any coincidence with
$+\ee_{a\inv h}$ only increases that coefficient. Hence
\eqref{eq:four-terms} is impossible. Thus $a=a\inv$ in $G_{2m}$,
which is exactly the last congruence in \eqref{eq:layer}.
Theorem~\ref{thm:kernel} now makes the symbol zero. The converse
implication to descent is immediate.
\end{proof}

The theorem is stronger than merely testing the normalized norm of the
new symbol. It detects a nonzero new-character component, so no
combination of old-field wedges can cancel it. This is the obstruction
that forces a modulus $2e$ in the final classification.

\subsection{Two finite-group difference tests}

For an odd conductor $m$, division by $2$ in symbol arguments is
well-defined. The following two tests will isolate the denominator and
numerator contributions of $J$.

\begin{lemma}[First difference]\label{lem:first}
For odd $m\ge1$ and $(a,m)=1$,
\begin{equation}\label{eq:first}
 \beta_m(a)-\beta_m(a/2)=0
 \quad\Longleftrightarrow\quad m\in\{1,3,5\}.
\end{equation}
\end{lemma}
\begin{proof}
The case $m=1$ is immediate. For $m\ge3$, put $x=\overline a$ and
$b=\overline2$ in $G_m$, and let $f_x=\ee_x-\ee_{x\inv}$.
By Theorem~\ref{thm:kernel}, the required equality is
$f_x=f_{x/b}$. If this vector is nonzero, its unique positive basis
coefficient identifies its index, so $x=x/b$ and $b=1$. If both
vectors are zero, $x^2=(x/b)^2=1$, so $b^2=1$. In either case
$b^2=1$, or $4\equiv\pm1\pmod m$, forcing odd $m$ to be $1$, $3$
or $5$. Conversely the groups at $3$ and $5$ have exponent at most
$2$, and all their symbols vanish.
\end{proof}

\begin{lemma}[Second difference]\label{lem:second}
For odd $m\ge1$ and $(a,m)=1$,
\begin{equation}\label{eq:second}
 \beta_m(2a)-2\beta_m(a)+\beta_m(a/2)=0
 \quad\Longleftrightarrow\quad a^2\equiv\pm1\pmod m.
\end{equation}
\end{lemma}
\begin{proof}
Again $m=1$ is automatic. In the group algebra of $G_m$, the image of
the left side is $(T_b+T_{b\inv}-2I)f_x$, where $T_b$ is translation
by $b=\overline2$. With the orthonormal basis $\ee_g$,
\[
 \big\langle f,(2I-T_b-T_{b\inv})f\big\rangle
 =\|f-T_bf\|^2.
\]
Vanishing of the second difference therefore forces $f_x=T_bf_x$.
If $f_x\ne0$, the unique positive coefficient of $f_x$ forces $b=1$.
For an odd conductor this means $m=1$ or $3$, at which $f_x=0$
already, a contradiction. Thus $f_x=0$, and the kernel theorem gives
the asserted congruence. Conversely $f_x=0$ makes its second difference
zero. The two positive and two negative translated terms agree with
the symbol expression even though a single translation of $f_x$ need
not itself have the form $f_y$.
\end{proof}
